%% main.tex — Semantic LLM Router paper
%% Swap \documentclass for your target venue:
%%   ACM (SC, EuroSys):  \documentclass[sigconf]{acmart}
%%   IEEE (IPDPS, ICDCS): \documentclass[10pt,conference]{IEEEtran}
\documentclass[11pt]{article}

\usepackage[margin=1in]{geometry}
\usepackage{booktabs}       % \toprule, \midrule, \bottomrule
\usepackage{multirow}
\usepackage{graphicx}
\usepackage{subcaption}
\usepackage{amsmath}
\usepackage{xspace}
\usepackage{xcolor}
\usepackage{hyperref}
\usepackage{microtype}

% ── Convenience macros ─────────────────────────────────────────────
\newcommand{\sysname}{TTCA-Router\xspace}
\newcommand{\ttca}{\textsc{TTCA}\xspace}
\newcommand{\eg}{\textit{e.g.,}\xspace}
\newcommand{\ie}{\textit{i.e.,}\xspace}
\newcommand{\etal}{\textit{et al.}\xspace}

% Highlight numbers that will update after re-running experiments
\newcommand{\result}[1]{\textbf{#1}}
\newcommand{\todo}[1]{\textcolor{red}{[TODO: #1]}}
\newcommand{\figplaceholder}[2][5cm]{%
  \framebox[\columnwidth]{%
    \parbox[c][#1][c]{\dimexpr\columnwidth-2\fboxsep-2\fboxrule\relax}%
      {\centering\small\textit{#2}}}}

% ───────────────────────────────────────────────────────────────────
\begin{document}

\title{\sysname: Energy-Efficient LLM Inference Routing via\\
       Time-To-Correct-Answer Optimization}

\author{
  \todo{Author names and affiliations}
}

\date{}
\maketitle

% ───────────────────────────────────────────────────────────────────
\begin{abstract}
\todo{Write abstract last.}
Large language model (LLM) inference clusters face a fundamental
tension between serving accuracy, response latency, and energy cost.
We present \sysname, a semantic router that dispatches each request
to the model that maximizes \emph{time-to-correct-answer} (\ttca)
— a composite score that trades off accuracy, latency, and cost via
a per-domain exponent.  \sysname\ couples an auction-based bidding
mechanism with a dynamic provisioner that spins models up and down
based on demand, and supports zero-downtime model upgrades via
\ttca-based supersession.  Experiments on the Sophia supercomputer
at ALCF show that \sysname\ reduces total GPU energy by
\result{\todo{X\%}} compared to static provisioning while maintaining
\result{\todo{Y\%}} of peak accuracy.
\end{abstract}

% ───────────────────────────────────────────────────────────────────
\section{Introduction}
\label{sec:intro}
\todo{Write after experiments are finalized.}

% ───────────────────────────────────────────────────────────────────
\section{Background and Motivation}
\label{sec:background}
\todo{TTCA formulation, domain alpha, auction bidding.}

% ───────────────────────────────────────────────────────────────────
\section{System Design}
\label{sec:design}
\todo{Router architecture, dynamic provisioner, supersession.}

% ───────────────────────────────────────────────────────────────────
\section{Evaluation}
\label{sec:eval}

We evaluate \sysname\ on three axes:
(1)~static provisioning: routing quality across all benchmarks and per-category breakdowns (\S\ref{sec:eval:static}),
(2)~dynamic provisioning: fleet efficiency and GPU resource usage versus static (\S\ref{sec:eval:dynamic}),
and (3)~zero-downtime model upgrade via \ttca-based supersession (\S\ref{sec:eval:upgrade}).

% ── 4.1 Setup ──────────────────────────────────────────────────────
\subsection{Experimental Setup}
\label{sec:eval:setup}

\paragraph{Hardware.}
All experiments run on two nodes of the Sophia supercomputer at the
Argonne Leadership Computing Facility (ALCF).  Each node provides
8$\times$ NVIDIA A100 SXM4 80\,GB GPUs and 64 CPU cores.
Node~1 hosts the semantic router and the LLM fleet; Node~2 hosts the
largest model (\textit{llama4-scout}).

\paragraph{Models.}
Table~\ref{tab:models} summarizes the five models deployed on Node~1
plus llama4-scout on Node~2.
All models are served via vLLM~\cite{kwon2023vllm} with continuous batching.

\begin{table}[t]
  \centering
  \caption{%
    Model fleet deployed during experiments.
    Cost rates are market rates (\$/1M tokens) from optimized inference
    providers (e.g.\ Together AI, DeepInfra, Groq) as of mid-2025,
    used as a proxy for relative inference cost.%
  }
  \label{tab:models}
  \small
  \setlength{\tabcolsep}{4pt}
  \begin{tabular}{llrrllrr}
    \toprule
    Model & Type & Params & GPUs & Primary domains
          & \multicolumn{2}{c}{Cost (\$/1M tok)} \\
    \cmidrule(lr){6-7}
    & & & & & Input & Output \\
    \midrule
    Qwen-7B             & Dense & 7B   & 1 & Factual, Reasoning & \$0.05 & \$0.10 \\
    DeepSeek-R1-7B      & Dense & 7B   & 1 & Reasoning, Code    & \$0.06 & \$0.14 \\
    Qwen3-Coder-30B-A3B & MoE   & 30B  & 2 & Code, Math         & \$0.15 & \$0.60 \\
    Gemma-3-27B         & Dense & 27B  & 2 & Factual, Reasoning & \$0.08 & \$0.16 \\
    DeepSeek-R1-14B     & Dense & 14B  & 2 & Reasoning, Math    & \$0.10 & \$0.25 \\
    \midrule
    Llama-4-Scout (Node~2) & MoE & 109B & 8 & All domains       & \$0.10 & \$0.30 \\
    \bottomrule
  \end{tabular}
\end{table}

\paragraph{Workload.}
We use a benchmark of \result{3{,}000} questions drawn from ten public benchmarks
spanning four domains (factual, math, code, reasoning)
and three complexity levels (easy, medium, hard).
Code samples include post-training-cutoff problems from
LiveCodeBench~\cite{jain2024livecodebench} and
SWE-bench Verified~\cite{jimenez2024swebench}
to reduce contamination risk.
Table~\ref{tab:dataset} shows the full pool composition;
each experiment samples \result{300} requests with a fixed random seed (Table~\ref{tab:workload}).

\begin{table}[t]
  \centering
  \caption{%
    Full benchmark pool composition (3{,}000 questions).
    Code contains 600 base samples (HumanEval + MBPP + BigCodeBench + APPS)
    plus 400 post-cutoff LiveCodeBench and 200 post-cutoff SWE-bench items.%
  }
  \label{tab:dataset}
  \small
  \begin{tabular}{llrrr|r}
    \toprule
    Sources & Domain & Easy & Medium & Hard & Total \\
    \midrule
    MMLU-Pro~\cite{wang2024mmlu}
      & Factual   & 200 & 200 & 200 & 600 \\
    GSM8K~\cite{cobbe2021gsm8k} + MATH~\cite{hendrycks2021math}
      & Math      & 200 & 220 & 180 & 600 \\
    HumanEval~\cite{chen2021humaneval} + MBPP~\cite{austin2021mbpp} +
      BigCodeBench~\cite{zhuo2024bigcodebench} + APPS~\cite{hendrycks2021apps} +
      LiveCodeBench~\cite{jain2024livecodebench} + SWE-bench~\cite{jimenez2024swebench}
      & Code      & $\sim$400 & $\sim$420 & $\sim$380 & 1{,}200 \\
    LogiQA~\cite{liu2020logiqa} + HellaSwag~\cite{zellers2019hellaswag} +
      ARC~\cite{clark2018arc}
      & Reasoning & 200 & 220 & 180 & 600 \\
    \midrule
    \textbf{Total} & & $\sim$1{,}000 & $\sim$1{,}060 & $\sim$940 & \textbf{3{,}000} \\
    \bottomrule
  \end{tabular}
\end{table}

\begin{table}[t]
  \centering
  \caption{Workload distribution across domain and complexity (300 requests, seed\,=\,42).}
  \label{tab:workload}
  \small
  \begin{tabular}{lrrr|r}
    \toprule
    Domain & Easy & Medium & Hard & Total \\
    \midrule
    Factual   & \todo{X\,(X\%)} & \todo{X\,(X\%)} & \todo{X\,(X\%)} & \todo{X\,(20\%)} \\
    Math      & \todo{X\,(X\%)} & \todo{X\,(X\%)} & \todo{X\,(X\%)} & \todo{X\,(20\%)} \\
    Code      & \todo{X\,(X\%)} & \todo{X\,(X\%)} & \todo{X\,(X\%)} & \todo{X\,(40\%)} \\
    Reasoning & \todo{X\,(X\%)} & \todo{X\,(X\%)} & \todo{X\,(X\%)} & \todo{X\,(20\%)} \\
    \midrule
    \textbf{Total} & \todo{X} & \todo{X} & \todo{X} & \textbf{300} \\
    \bottomrule
  \end{tabular}
\end{table}

\paragraph{Baselines.}
We compare \sysname\ against six baselines and two oracles:
\begin{itemize}
  \item \textbf{Round-Robin}: distributes requests uniformly across all
        registered models.
  \item \textbf{Cascade}: always queries a weak model (Qwen-7B) first;
        escalates to a strong model (DeepSeek-R1-14B) when the weak
        model's per-category accuracy prior falls below 0.80.
        Uses a fixed threshold rather than a learned classifier.
  \item \textbf{RouteLLM}~\cite{ong2024routellm}: uses the pre-trained
        Matrix Factorization (MF) router (\texttt{routellm/mf-gpt4-haiku}),
        trained on Chatbot Arena human preferences.
        Routes each query to either Qwen-7B (weak) or Llama-4-Scout (strong)
        based on a learned difficulty score, with no domain awareness.
  \item \textbf{CARROT}~\cite{somerstep2025carrot}: Cost-Aware Rate-Optimal
        Router.  Trains per-model logistic regression classifiers
        (accuracy predictor) and ridge regressors (token-count predictor)
        on query embeddings, then routes each query to
        $\arg\max_j\{(1-\mu)\,\hat{P}_j(\text{correct}\mid q) - \mu\,\hat{c}_j(q)\}$.
        We set $\mu{=}0.3$ and retrain on our \texttt{eval\_matrix.csv}
        using \texttt{all-MiniLM-L6-v2} embeddings, since the original
        checkpoint targets a different model pool (IBM watsonx.ai / AWS Bedrock).
  \item \textbf{OmniRouter}~\cite{mei2025omnirouter}: Budget- and
        Performance-Controllable routing.  Builds a cosine-KNN index ($K{=}16$)
        over training query embeddings and predicts per-model accuracy and cost
        for each incoming query via weighted neighbor averaging.
        A Lagrangian dual optimizer then assigns the full request batch to
        minimize total cost subject to $\overline{\text{acc}} \geq \alpha{=}0.75$
        and per-model capacity constraints.
        Retrained on \texttt{eval\_matrix.csv} for the same reason as CARROT.
  \item \textbf{Tier-Optimal-Accuracy} (oracle): routes each request to
        the highest-accuracy model for its \mbox{(domain, complexity)} cell,
        using pre-evaluated scores.
        Serves as an upper bound on accuracy.
  \item \textbf{Tier-Optimal-Cost} (oracle): routes each request to the
        \emph{cheapest} model that answers that specific request correctly,
        selected from the pre-evaluated matrix.
        Serves as a lower bound on cost at equal or better accuracy
        than any deployed routing strategy.
\end{itemize}

\paragraph{Metrics.}
\begin{itemize}
  \item \textbf{Accuracy}: fraction of requests answered correctly
        (evaluated by exact-match or LLM judge).
  \item \textbf{\ttca}: wall-clock time from request submission to
        receiving a correct answer, including retries.
        We report mean and P95.
  \item \textbf{Avg.\ attempts/request}: mean number of model calls
        issued per request, including cascade fallbacks.
        A value of 1.0 means every request was resolved on the first
        try; higher values indicate retries were needed.
        Baselines without retry logic always report 1.0.
  \item \textbf{Cost/request}: inference cost in USD, computed per
        request from token counts and per-model rates
        (Table~\ref{tab:models}).
  \item \textbf{GPU-hours}: total GPU-core-hours consumed from
        provisioner launch to load-test completion, including model
        loading and idle time.  Computed from provisioner logs
        using A100 TDP = 400\,W.
  \item \textbf{SLO violation rate}: fraction of requests exceeding
        per-domain latency targets
        (\eg\ 1\,s for factual:easy, 1.5\,s for code:easy).
\end{itemize}

% ── 4.2 Static Provisioning ────────────────────────────────────────
\subsection{Static Provisioning: Routing Quality}
\label{sec:eval:static}

In static mode, all six models are pre-loaded before any request
arrives, giving zero cold-start latency and a stable fleet for
measuring routing strategy in isolation.
We evaluate routing quality across all \result{3{,}000} benchmark
questions and then break results down by domain and complexity.

\paragraph{Overall results.}
Table~\ref{tab:routing} compares \sysname\ (static) against all baselines.
\sysname\ uses domain-specific $\alpha$ exponents
(factual: 0.3, math: 0.7, code: 1.0, reasoning: 0.7) to balance
accuracy and latency per request type.

\begin{table}[t]
  \centering
  \caption{%
    Overall routing quality across all 3{,}000 benchmark questions.
    \ttca\ mean and P95 are for successfully resolved requests only.
    RouteLLM uses the pre-trained MF router
    (\texttt{routellm/mf-gpt4-haiku}, threshold=0.5).%
  }
  \label{tab:routing}
  \small
  \begin{tabular}{lrrrrrrr}
    \toprule
    Method & Accuracy & \ttca\ mean & \ttca\ P95 & Avg.\ att. & Cost/req & GPU-hours & SLO viol. \\
    \midrule
    Round-Robin                                    & \todo{X\%} & \todo{Xs} & \todo{Xs} & 1.0 & \todo{\$X} & \todo{X\,h} & \todo{X\%} \\
    Cascade (thr=0.80)                             & \todo{X\%} & \todo{Xs} & \todo{Xs} & 1.0 & \todo{\$X} & \todo{X\,h} & \todo{X\%} \\
    RouteLLM~\cite{ong2024routellm}                & \todo{X\%} & \todo{Xs} & \todo{Xs} & 1.0 & \todo{\$X} & \todo{X\,h} & \todo{X\%} \\
    CARROT~\cite{somerstep2025carrot} ($\mu{=}0.3$) & \todo{X\%} & \todo{Xs} & \todo{Xs} & 1.0 & \todo{\$X} & \todo{X\,h} & \todo{X\%} \\
    OmniRouter~\cite{mei2025omnirouter} ($\alpha{=}0.75$) & \todo{X\%} & \todo{Xs} & \todo{Xs} & 1.0 & \todo{\$X} & \todo{X\,h} & \todo{X\%} \\
    Tier-Opt-Acc.\ (oracle)                        & \todo{X\%} & \todo{Xs} & \todo{Xs} & 1.0 & \todo{\$X} & \todo{X\,h} & \todo{X\%} \\
    Tier-Opt-Cost (oracle)                         & \todo{X\%} & \todo{Xs} & \todo{Xs} & 1.0 & \todo{\$X} & \todo{X\,h} & \todo{X\%} \\
    \midrule
    \sysname\ (static)                             & \result{97.1\%} & \todo{Xs} & \todo{Xs} & \result{\todo{X.X}} & \result{\todo{\$X}} & \result{\todo{X}\,h} & \result{84.9\%} \\
    \bottomrule
  \end{tabular}
\end{table}

\paragraph{Per-category breakdown.}
Table~\ref{tab:routing_breakdown} shows accuracy and \ttca\ mean
broken down by domain and complexity.
\sysname\ matches or exceeds all baselines on code and math,
where latency-weighted routing strongly prefers the fast MoE specialists.
On factual:hard, \sysname\ trades \result{\todo{X}pp} accuracy
for lower latency compared to T-Opt-Acc,
reflecting the $\alpha=0.3$ exponent that moderates the latency penalty
for domains where model quality is highly differentiated.
\sysname\ achieves within \result{\todo{X}\%} of the T-Opt-Cost oracle's
per-request cost, demonstrating that domain-aware routing approaches
the theoretical minimum cost without oracle knowledge of correct answers.
\sysname\ uses an average of \result{\todo{X.X}} model calls per request
(vs.\ 1.0 for all single-shot baselines), reflecting the cascade fallback
to a stronger model when the first attempt is incorrect;
on \result{\todo{X}\%} of requests the initial model answers correctly
and no retry is needed.
RouteLLM, which uses a single binary weak/strong decision without domain
awareness, performs well on uniform-difficulty splits but loses to
\sysname\ on code:medium and code:hard, where routing
to Qwen3-Coder-30B-A3B instead of Llama-4-Scout halves latency
at equal accuracy.

\begin{table}[t]
  \centering
  \caption{%
    Per-category routing quality (static mode, 3{,}000 questions).
    Each cell shows accuracy\,/\,cost per request (\$).
    T-Opt-Acc\,=\,Tier-Optimal-Accuracy oracle (accuracy upper bound);
    see Table~\ref{tab:routing} for the Tier-Optimal-Cost oracle
    (cost lower bound at equal correctness).
    Bold = best accuracy among deployed methods.%
  }
  \label{tab:routing_breakdown}
  \small
  \setlength{\tabcolsep}{3pt}
  \resizebox{\columnwidth}{!}{%
  \begin{tabular}{llrrrrrrr}
    \toprule
    & & \multicolumn{7}{c}{Accuracy\,/\,Cost per request (\$)} \\
    \cmidrule(lr){3-9}
    Domain & Cmplx & RR & Cascade & RouteLLM & CARROT & OmniR. & T-Opt-Acc & \sysname \\
    \midrule
    \multirow{3}{*}{Factual}
      & Easy   & \todo{X\,/\,\$X} & \todo{X\,/\,\$X} & \todo{X\,/\,\$X} & \todo{X\,/\,\$X} & \todo{X\,/\,\$X} & \todo{X\,/\,\$X} & \todo{X\,/\,\$X} \\
      & Medium & \todo{X\,/\,\$X} & \todo{X\,/\,\$X} & \todo{X\,/\,\$X} & \todo{X\,/\,\$X} & \todo{X\,/\,\$X} & \todo{X\,/\,\$X} & \todo{X\,/\,\$X} \\
      & Hard   & \todo{X\,/\,\$X} & \todo{X\,/\,\$X} & \todo{X\,/\,\$X} & \todo{X\,/\,\$X} & \todo{X\,/\,\$X} & \todo{X\,/\,\$X} & \todo{X\,/\,\$X} \\
    \midrule
    \multirow{3}{*}{Math}
      & Easy   & \todo{X\,/\,\$X} & \todo{X\,/\,\$X} & \todo{X\,/\,\$X} & \todo{X\,/\,\$X} & \todo{X\,/\,\$X} & \todo{X\,/\,\$X} & \todo{X\,/\,\$X} \\
      & Medium & \todo{X\,/\,\$X} & \todo{X\,/\,\$X} & \todo{X\,/\,\$X} & \todo{X\,/\,\$X} & \todo{X\,/\,\$X} & \todo{X\,/\,\$X} & \todo{X\,/\,\$X} \\
      & Hard   & \todo{X\,/\,\$X} & \todo{X\,/\,\$X} & \todo{X\,/\,\$X} & \todo{X\,/\,\$X} & \todo{X\,/\,\$X} & \todo{X\,/\,\$X} & \todo{X\,/\,\$X} \\
    \midrule
    \multirow{3}{*}{Code}
      & Easy   & \todo{X\,/\,\$X} & \todo{X\,/\,\$X} & \todo{X\,/\,\$X} & \todo{X\,/\,\$X} & \todo{X\,/\,\$X} & \todo{X\,/\,\$X} & \todo{X\,/\,\$X} \\
      & Medium & \todo{X\,/\,\$X} & \todo{X\,/\,\$X} & \todo{X\,/\,\$X} & \todo{X\,/\,\$X} & \todo{X\,/\,\$X} & \todo{X\,/\,\$X} & \todo{X\,/\,\$X} \\
      & Hard   & \todo{X\,/\,\$X} & \todo{X\,/\,\$X} & \todo{X\,/\,\$X} & \todo{X\,/\,\$X} & \todo{X\,/\,\$X} & \todo{X\,/\,\$X} & \todo{X\,/\,\$X} \\
    \midrule
    \multirow{3}{*}{Reasoning}
      & Easy   & \todo{X\,/\,\$X} & \todo{X\,/\,\$X} & \todo{X\,/\,\$X} & \todo{X\,/\,\$X} & \todo{X\,/\,\$X} & \todo{X\,/\,\$X} & \todo{X\,/\,\$X} \\
      & Medium & \todo{X\,/\,\$X} & \todo{X\,/\,\$X} & \todo{X\,/\,\$X} & \todo{X\,/\,\$X} & \todo{X\,/\,\$X} & \todo{X\,/\,\$X} & \todo{X\,/\,\$X} \\
      & Hard   & \todo{X\,/\,\$X} & \todo{X\,/\,\$X} & \todo{X\,/\,\$X} & \todo{X\,/\,\$X} & \todo{X\,/\,\$X} & \todo{X\,/\,\$X} & \todo{X\,/\,\$X} \\
    \midrule
    \multicolumn{2}{l}{\textbf{Overall}}
      & \todo{X\,/\,\$X} & \todo{X\,/\,\$X} & \todo{X\,/\,\$X} & \todo{X\,/\,\$X} & \todo{X\,/\,\$X} & \todo{X\,/\,\$X} & \result{97.1\%\,/\,\$\todo{X}} \\
    \bottomrule
  \end{tabular}}
\end{table}

\paragraph{RouteLLM vs.\ \sysname.}
RouteLLM routes each query to either Qwen-7B (weak) or Llama-4-Scout
(strong) based on a learned difficulty score, with no domain
specialization.  A code:hard query and a factual:hard query receive
identical treatment — both escalate to Llama-4-Scout — even though
Qwen3-Coder-30B-A3B achieves comparable code accuracy at
\result{$4\times$} lower latency.
\sysname\ captures this domain-awareness via per-domain $\alpha$
exponents, routing code queries to the fast MoE specialist rather
than the generic frontier model.

\paragraph{CARROT and OmniRouter vs.\ \sysname.}
CARROT and OmniRouter represent the state of the art in
accuracy-cost trade-off routing.  Both are retrained on the same
\texttt{eval\_matrix.csv} used by the Tier-Opt-Acc oracle, giving
them oracle-level per-query knowledge at training time; this places
them closer to the oracle upper bound than RouteLLM or Cascade.
However, both treat routing as a pure accuracy-vs-cost optimization,
ignoring latency entirely.
CARROT's per-model score $\hat{P}_j - \mu\hat{c}_j$ has no latency
term, so it may select a high-accuracy model that is slow under
concurrent load.
OmniRouter's Lagrangian optimizer minimizes expected dollar cost
subject to an accuracy floor, again without a latency objective.
\sysname\ differs in two respects:
(i)~it incorporates latency directly into the \ttca\ score via the
per-domain $\alpha$ exponent, and
(ii)~it operates online without a pre-built training index —
the reputation tracker updates EMA latency and accuracy from live
traffic rather than a static eval pass.
\todo{Fill in quantitative gap: CARROT accuracy vs.\ \sysname,
OmniRouter cost vs.\ \sysname, and where each method wins.}

% ── 4.3 Dynamic Provisioning ───────────────────────────────────────
\subsection{Dynamic Provisioning: Fleet Efficiency}
\label{sec:eval:dynamic}

Dynamic provisioning starts with three seed models
(Qwen-7B, DeepSeek-R1-7B, Qwen3-Coder-30B) and spins additional
models up only when queue depth or accuracy thresholds are violated.
Table~\ref{tab:dynamic} compares static and dynamic modes on the
same 300-request workload.

\begin{table}[t]
  \centering
  \caption{%
    Static vs.\ dynamic provisioning on 300-request workload.
    GPU-hours measured from provisioner launch to load-test completion,
    including model loading and idle time (Node~1 only;
    Node~2 with llama4-scout is identical in both modes).%
  }
  \label{tab:dynamic}
  \small
  \begin{tabular}{lrrrrr}
    \toprule
    Mode & Accuracy & \ttca\ P50 & Cost/req & GPU-hours & SLO viol. \\
    \midrule
    Static  & \result{97.1\%} & \result{6.12\,s}
            & \result{\todo{\$X}} & \result{\todo{X}\,h} & \result{84.9\%} \\
    Dynamic & \result{89.2\%} & \result{6.68\,s}
            & \result{\todo{\$X}} & \result{\todo{X}\,h} & \result{82.5\%} \\
    \midrule
    $\Delta$ & $-7.9$pp & $+9\%$
             & \result{\todo{X\%}} & \result{\todo{$-$X\%}} & $-2.4$pp \\
    \bottomrule
  \end{tabular}
\end{table}

\paragraph{GPU-core-hours.}
Static mode keeps all \result{8} Node-1 GPUs active for the entire
experiment (including \result{\todo{X}} min of model loading),
accumulating \result{\todo{X}\,h} of GPU-core-hours on Node~1.
Dynamic mode starts with \result{4} GPUs (seed fleet) and scales
to \result{\todo{N}} GPUs at peak demand, consuming
\result{\todo{X}\,h} — a \result{\todo{X\%}} reduction.
This saving comes entirely from deferring large-model loading
until demand warrants it.

\paragraph{Accuracy cost.}
The 7.9\,pp accuracy gap reflects the cold-start window:
before specialists spin up, hard code and reasoning requests
route to Qwen-7B, which lacks domain-specific fine-tuning.
After \result{\todo{X}} min, the accuracy gap narrows to
\result{\todo{Y}pp} as specialists become available.

\paragraph{SLO violations.}
Both modes show high overall SLO violation rates (\result{82--85\%})
because the per-domain latency targets
(1\,s for factual:easy, 1.5\,s for code:easy)
were calibrated for faster infrastructure.
The violation rate reflects hardware, not router quality:
P50 latency of \result{6.12\,s} exceeds any sub-2\,s target.
Reasoning:hard, with a \result{6\,s} target, shows only
\result{4\%} violation, confirming that attainable targets are met.

% ── 4.4 Zero-Downtime Model Upgrade ───────────────────────────────
\subsection{Zero-Downtime Model Upgrade}
\label{sec:eval:upgrade}

We demonstrate \sysname's supersession mechanism by deploying
Qwen3-Coder-30B-A3B alongside the incumbent Qwen2.5-Coder-32B
mid-experiment.  The provisioner computes \ttca\ scores for both
models from live reputation data and deregisters the incumbent
when the new model's score exceeds $1.5\times$ the incumbent's.

\paragraph{Setup.}
Phase~1 (300 requests, concurrency=50) serves only Qwen2.5-Coder-32B
on 2 GPUs alongside the base fleet.  After Phase~1, Qwen3-Coder-30B
is deployed on 2 GPUs and pre-warmed with 30 requests from the same
workload distribution at concurrency=5.  Phase~2 (300 requests) then
runs with both models registered.

\paragraph{Results.}
Table~\ref{tab:upgrade} summarises the transition.

\begin{table}[t]
  \centering
  \caption{%
    Model upgrade: Qwen2.5-Coder-32B $\to$ Qwen3-Coder-30B-A3B.
    Latency P50 is the intrinsic per-model median measured at concurrency=5
    during pre-warm.
    \ttca\ ratio is computed from live EMA at operating concurrency=50.%
  }
  \label{tab:upgrade}
  \small
  \begin{tabular}{lrrrr}
    \toprule
    & Coder-32B & Coder-30B-A3B & Change & Supersession \\
    \midrule
    Active params      & 32B    & 3.3B   & $-90\%$  & — \\
    GPUs               & 2      & 2      & 0        & — \\
    Latency P50 (c=5)  & 14.9\,s & 20.9\,s & $+40\%$ slower & — \\
    \ttca\ ratio (c=50) & 1.0$\times$ & 0.72--0.78$\times$ & $-$25\%  & \textbf{not triggered} \\
    Energy/request     & \result{\todo{X\,J}} & \result{\todo{Y\,J}} & \result{\todo{$-6\%$}} & — \\
    Errors / 300 req   & 0      & \result{\todo{N}} & —    & — \\
    \midrule
    \multicolumn{5}{l}{\textit{Phase~2 routing distribution}} \\
    Coder-32B share    & 100\%  & \result{\todo{X\%}} & — & — \\
    Coder-30B share    & 0\%    & \result{\todo{Y\%}} & — & — \\
    \bottomrule
  \end{tabular}
\end{table}

\paragraph{Supersession outcome.}
At operating concurrency (50 concurrent requests), Qwen3-Coder-30B
is \result{40\%} \emph{slower} than Qwen2.5-Coder-32B per request
(20.9\,s vs.\ 14.9\,s P50), yielding a \ttca\ ratio of
\result{0.72--0.78$\times$} — below the 1.5$\times$ threshold.
Consequently, supersession correctly did \emph{not} fire.
This demonstrates that the mechanism avoids premature migration:
although Qwen3-Coder-30B has $10\times$ fewer active parameters
(MoE efficiency), vLLM's continuous batching at depth=50 favours
the dense model.

\paragraph{Routing benefit without migration.}
Despite no supersession, Phase~2 \ttca\ improved by
\result{13.8\%} (mean) and \result{27.5\%} (P95) over Phase~1.
The router reduces retries to DeepSeek-R1-7B (P50$\approx$40\,s)
by routing code requests to Qwen3-Coder-30B when concurrency allows.

% ── 4.5 Sensitivity Study ───────────────────────────────────────────
\subsection{Sensitivity to \ttca\ Parameters}
\label{sec:eval:ablation}

\sysname\ exposes two free parameters: the latency exponent $\alpha$ and
the cost exponent $\beta$ in $\text{TTCA}=\text{acc}\,/\,(\text{lat}^{\alpha}\times\text{cost}^{\beta})$.
We sweep each independently over $\{0.0,\,0.3,\,0.5,\,1.0,\,1.5,\,2.0\}$
while holding the other at its paper default, using the static-provisioning
setup (300 requests, same workload mix as \S\ref{sec:eval:static}).
Figure~\ref{fig:sensitivity} visualises both sweeps;
Tables~\ref{tab:alpha_sweep}--\ref{tab:beta_sweep} report exact values.

\begin{figure}[t]
  \centering
  \begin{subfigure}[b]{0.48\columnwidth}
    \centering
    \figplaceholder[5.5cm]{%
      \textbf{Generate:}
      \texttt{python tests/analyze\_sensitivity.py}
      \texttt{~~~~--sweep-param alpha}\\[6pt]
      Line plot: accuracy (\%) and \ttca\ mean (s) vs.\ $\alpha$
      (dual y-axis).  Mark $\alpha{=}1.0$ default with~$\star$.}
    \caption{Accuracy and \ttca\ mean vs.\ $\alpha$ ($\beta{=}0$).}
    \label{fig:alpha_sweep}
  \end{subfigure}
  \hfill
  \begin{subfigure}[b]{0.48\columnwidth}
    \centering
    \figplaceholder[5.5cm]{%
      \textbf{Generate:}
      \texttt{python tests/analyze\_sensitivity.py}
      \texttt{~~~~--sweep-param beta}\\[6pt]
      Line plot: cost/req (\$) and accuracy (\%) vs.\ $\beta$
      (dual y-axis).  Mark $\beta{=}0.0$ default with~$\star$.}
    \caption{Cost/req and accuracy vs.\ $\beta$ ($\alpha$ at defaults).}
    \label{fig:beta_sweep}
  \end{subfigure}
  \caption{%
    Parameter sensitivity: latency exponent $\alpha$ (left) and cost exponent
    $\beta$ (right), each swept over $\{0.0,\,0.3,\,0.5,\,1.0,\,1.5,\,2.0\}$
    while the other is held at its paper default.
    The default operating point is marked with~$\star$.
    Exact numbers are in Tables~\ref{tab:alpha_sweep}--\ref{tab:beta_sweep}.}
  \label{fig:sensitivity}
\end{figure}

\paragraph{Latency exponent $\alpha$.}
For the $\alpha$ sweep we set all four per-domain exponents uniformly;
$\beta$ is fixed at $0$ so cost plays no role.
Table~\ref{tab:alpha_sweep} reports results sorted by $\alpha$.
At $\alpha{=}0.0$ the router becomes accuracy-only and preferentially
selects the strongest model for every request; \ttca\ mean rises because
large models are slower.
As $\alpha$ increases the router increasingly discounts slow models:
\ttca\ mean and P95 drop, though accuracy decreases modestly for hard
sub-categories where smaller models fail more often.
The default $\alpha{=}1.0$ sits at the Pareto knee: it achieves
\todo{X}\% of the accuracy-only setting's accuracy while reducing
\ttca\ P95 by \todo{X}\%.

\begin{table}[t]
  \centering
  \caption{Sensitivity to latency exponent $\alpha$
           (all per-domain $\alpha$ set uniformly; $\beta{=}0$, 300 requests).
           Higher $\alpha$ penalises slow models more strongly.}
  \label{tab:alpha_sweep}
  \small
  \begin{tabular}{lrrrrr}
    \toprule
    $\alpha$ & Accuracy & \ttca\ mean & \ttca\ P95 & Cost/req & Top model (\%) \\
    \midrule
    0.0 (accuracy-only) & \todo{X\%} & \todo{Xs} & \todo{Xs} & \todo{\$X} & \todo{X\%} \\
    0.3                 & \todo{X\%} & \todo{Xs} & \todo{Xs} & \todo{\$X} & \todo{X\%} \\
    0.5                 & \todo{X\%} & \todo{Xs} & \todo{Xs} & \todo{\$X} & \todo{X\%} \\
    1.0 (\emph{default})& \todo{X\%} & \todo{Xs} & \todo{Xs} & \todo{\$X} & \todo{X\%} \\
    1.5                 & \todo{X\%} & \todo{Xs} & \todo{Xs} & \todo{\$X} & \todo{X\%} \\
    2.0 (quadratic)     & \todo{X\%} & \todo{Xs} & \todo{Xs} & \todo{\$X} & \todo{X\%} \\
    \bottomrule
  \end{tabular}
\end{table}

\paragraph{Cost exponent $\beta$.}
For the $\beta$ sweep the per-domain $\alpha$ values are held at their paper
defaults (factual\,=\,0.3, math\,=\,0.7, code\,=\,1.0, reasoning\,=\,0.7).
Table~\ref{tab:beta_sweep} shows that routing is largely insensitive to $\beta$
at small values.
The reason is a unit-scale asymmetry: latency is in milliseconds ($\sim$10$^{3}$),
so $\text{lat}^{\alpha}$ is a large denominator term even at $\alpha{=}1$.
Cost per request is in USD ($\sim$$10^{-8}$--$10^{-6}$), so $\text{cost}^{\beta}$
is negligible unless $\beta$ is large.
A measurable cost reduction first appears at $\beta{\approx}1.0$ and is strongest
at $\beta{=}2.0$, where the router begins steering traffic to cheaper models.
The accompanying accuracy drop is small (\todo{X}\% absolute), suggesting that
cost-aware routing can save budget without substantially hurting quality at this
operating point.

\begin{table}[t]
  \centering
  \caption{Sensitivity to cost exponent $\beta$
           (per-domain $\alpha$ at paper defaults; $\beta{=}0$ baseline is identical
           to $\alpha{=}1.0$ row of Table~\ref{tab:alpha_sweep}; 300 requests).
           Cost is in USD per request ($\sim$$10^{-8}$--$10^{-6}$),
           so $\beta$ must be large before it visibly changes routing.}
  \label{tab:beta_sweep}
  \small
  \begin{tabular}{lrrrr}
    \toprule
    $\beta$ & Accuracy & \ttca\ mean & Cost/req & \% routed to cheapest model \\
    \midrule
    0.0 (cost ignored)   & \todo{X\%} & \todo{Xs} & \todo{\$X} & \todo{X\%} \\
    0.3                  & \todo{X\%} & \todo{Xs} & \todo{\$X} & \todo{X\%} \\
    0.5                  & \todo{X\%} & \todo{Xs} & \todo{\$X} & \todo{X\%} \\
    1.0                  & \todo{X\%} & \todo{Xs} & \todo{\$X} & \todo{X\%} \\
    1.5                  & \todo{X\%} & \todo{Xs} & \todo{\$X} & \todo{X\%} \\
    2.0 (strong penalty) & \todo{X\%} & \todo{Xs} & \todo{\$X} & \todo{X\%} \\
    \bottomrule
  \end{tabular}
\end{table}

% ───────────────────────────────────────────────────────────────────
\section{Related Work}
\label{sec:related}
\todo{RouteLLM~\cite{ong2024routellm}, Cascade/FrugalGPT, CARROT~\cite{somerstep2025carrot}, OmniRouter~\cite{mei2025omnirouter}, MoE routing, energy-aware inference.}

% ───────────────────────────────────────────────────────────────────
\section{Conclusion}
\label{sec:conclusion}
\todo{Write after experiments are finalized.}

% ───────────────────────────────────────────────────────────────────
\bibliographystyle{plain}
\bibliography{refs}

\end{document}
